% Revised Section 5 — Experimental Results
%
% Based on manuscript_draft.pdf pages 9–11.
% All numbers verified against:
%   results/contribution1/primary_calibration.csv
%   results/contribution1/sensitivity_sweep.csv
%   results/contribution1/uniform_comparison.csv
%   results/contribution1/scorecard.csv
%   results/auroc_per_class.csv
%   results/coverage_simulation.csv
%
% Changes from the submitted draft:
%   5.2.1 — "nonconformity scores drawn from Beta(7,3) and Beta(3,7)" corrected
%            to "probability scores"; the LTT tests f_theta(x)_k >= lambda
%            directly, never inverts to 1 - score.
%   5.2.2 — Rewrote the "Two of the uncertifiable Critical/Important results"
%            paragraph. The original claimed these classes were uncertifiable
%            due to sample size "rather than test FNR" — implying other
%            uncertifiable classes were due to FNR, which is wrong: all
%            uncertifiable results arise from the LTT certification floor not
%            being met. The revised paragraph explains the two-floor distinction
%            (quantile-existence vs. LTT) and covers all affected classes
%            (AFIB, _AVB, CLBBB, ISCA, ISCI), not just the Critical-tier pair.
%   5.2.4 — "effective FNR bound of 0.054" corrected to "FNR guarantee of 0.10".
%            0.054 is lambda_c (the threshold), not the FNR bound; the bound
%            certified by the binomial LTT at alpha=0.10 is 0.10.
%   5.1   — Section header only (placeholder); unchanged.
%
% Numbers verified correct and unchanged:
%   MC simulation: violation=0.078, mean lambda_c=0.497, mean FNR=0.088
%   Hoeffding grid-Bonferroni: violation=0.000, lambda_c=0.418, FNR=0.033
%   Hoeffding monotone: violation=0.000, lambda_c=0.473, FNR=0.067
%   Certified 11, trivial 2, uncertifiable 11
%   AMI: lambda_c=0.0060, test FNR=0.0131, flag=54.0%
%   IMI: lambda_c=0.0020, test FNR=0.0061, flag=76.1%
%   Benign: 9/14 certified, FNR range 0.0000 (NST_) – 0.0901 (STTC)
%   Flag range: 9.8% (STACH) – 96.1% (IVCD)
%   14/24 violations for both uncorrected baselines
%   Safety failures: 6 (AFIB, AMI, IMI, _AVB, ISCA, ISC_)
%   Efficiency losses: 9 (IRBBB, IVCD, LAFB/LPFB, LVH, NST_, SARRH, SBRAD, STACH, STTC)
%   AUROC range: 0.746 (IVCD) – 0.998 (CRBBB); AFIB=0.983, CLBBB=0.997
%   34 certified (class, config) pairs across 16 distinct classes

% ─────────────────────────────────────────────────────────────────────────────
\section{Experimental Results}

\subsection{Experimental Setup}
% [Placeholder — content to be added]

\subsection{Acuity-Weighted Class-Conditional Risk Control: Results}

\subsubsection{Validity of the testing procedure}

Certification used the Learn-Then-Test (LTT) framework with a one-sided binomial
risk test and a monotone-boundary search over the threshold $\lambda$, controlling
the family-wise error rate at $\delta{=}0.10$ across the 24 supported classes via
Bonferroni correction (per-class $\delta\approx 0.0042$).
The validity of this configuration was assessed by Monte Carlo simulation
($n_\mathrm{trials}{=}1000$, $n_\mathrm{cal}{=}2000$, target $\alpha{=}0.10$,
$\delta{=}0.10$, positive and negative probability scores drawn from
% CHANGED: "nonconformity scores" -> "probability scores".
% The LTT procedure tests f_theta(x)_k >= lambda directly.  Nothing is inverted
% to 1-score; calling them nonconformity scores is incorrect in this context.
$\mathrm{Beta}(7,3)$ and $\mathrm{Beta}(3,7)$ respectively, with the true
false-negative rate evaluated analytically from the Beta CDF).
The adopted binomial monotone-boundary variant produced an empirical violation rate
of 0.078, within the permitted bound $\delta{=}0.10$, at a mean certified threshold
$\bar\lambda_c{=}0.497$ and mean FNR 0.088 (Table~\ref{tab:coverage_sim}).
Two alternative estimators were evaluated for reference: a Hoeffding bound with
grid-Bonferroni correction (violation~$={=}$~0.000, $\bar\lambda_c{=}0.418$, mean
FNR~0.033) and a Hoeffding bound with monotone boundary
(violation~$={=}$~0.000, $\bar\lambda_c{=}0.473$, mean FNR~0.067).
All three configurations satisfied the coverage requirement.

\subsubsection{Primary calibration outcomes}

Across the 24 supported classes, the acuity-weighted procedure certified a
class-conditional FNR guarantee at the tier-specific target for 11~classes,
returned a trivial threshold ($\lambda{=}0$, indicating no discriminative
selectivity) for 2~classes, and returned no certificate for 11~classes
(Table~\ref{tab:primary_calibration}).
For every certified class, the FNR measured on the held-out test fold (fold~10)
fell at or below the class's target $\alpha$.

In the Critical tier ($\alpha{=}0.02$), AMI and IMI were certified, with test FNRs
of 0.0131 and 0.0061 at certified thresholds $\lambda_c{=}0.0060$ and 0.0020 and
corresponding flag rates of 54.0\% and 76.1\%.
The remaining three Critical-tier classes (AFIB, LMI, \texttt{\_AVB}) were
uncertifiable.
In the Important tier ($\alpha{=}0.05$), no class was certified: four classes were
uncertifiable and ISC\_ returned a trivial threshold.
In the Benign tier ($\alpha{=}0.10$), 9 of 14 classes were certified, with test
FNRs ranging from 0.0000 (NST\_) to 0.0901 (STTC) and flag rates ranging from
9.8\% (STACH) to 96.1\% (IVCD).

% CHANGED: rewrote this paragraph entirely.
% Original: "Two of the uncertifiable Critical/Important results were associated
%            with calibration-fold sample size rather than with test FNR."
% Problem:  ALL uncertifiable results are due to the LTT certification floor not
%           being met -- this is a function of n_k, alpha_k, and delta_k alone,
%           independent of the actual score distribution.  Framing two classes as
%           sample-size-limited "rather than FNR-limited" implies the others are
%           FNR-limited, which is wrong.  The correct distinction is between
%           classes that fall between the two floors (quantile-existence vs. LTT)
%           and those below both floors.  CLBBB and ISCA are in the same
%           between-floors situation as AFIB and _AVB; the original ignored them.
All uncertifiable results arise from the LTT certification floor not being met
at $\delta_k\approx 0.0042$: the condition $(1-\alpha_k)^{n_k^+}\le\delta_k$
depends only on sample count, $\alpha_k$, and the Bonferroni-corrected
$\delta_k$---not on the test-set score distribution.
The notable split is between classes that fall \emph{between} the two floors
(quantile-existence floor met, LTT floor not) and those below both.
In the Critical tier, AFIB ($n{=}151$) and \texttt{\_AVB} ($n{=}83$) both exceed
the quantile-existence floor ($n\ge 49$) and therefore appear as Supported in
Table~\ref{tab:sample_counts}, but neither reaches the LTT certification floor
of~272; LMI ($n{=}20$) falls below the quantile-existence floor.
In the Important tier, CLBBB ($n{=}54$) and ISCA ($n{=}92$) are in the same
between-floors position (quantile-existence floor $n\ge 19$ met; LTT floor
$n\ge 107$ not); ISCI ($n{=}39$) exceeds the quantile-existence floor marginally
but also falls short of the LTT floor; SVARR ($n{=}15$) is below both floors.
The two trivial classes (ISC\_, CRBBB) form a distinct category: both meet the LTT
certification floor (ISC\_ $n{=}125\ge 107$; CRBBB $n{=}55\ge 53$), but their
positive-class scores are concentrated near zero, leaving no threshold above
$\lambda{=}0$ that passes the binomial test.

\subsubsection{Comparison against uncorrected calibration baselines}

To assess whether per-class miscoverage control requires the statistical correction
applied by the binomial LTT procedure, two baselines without such correction were
evaluated: temperature-scaled thresholding~\cite{} and MC Dropout
thresholding~\cite{}.
Each baseline selects, per class, the threshold $\lambda$ at which
validation-fold FNR first falls at or below $\alpha_\mathrm{tier}$, with no
adjustment for sampling variability between the validation and test folds.

Across the 24 supported classes, the binomial LTT procedure produced no violations
of the target $\alpha_\mathrm{tier}$ on the test fold (0/24, 0\%).
Both uncorrected baselines violated their target on 14 of 24 classes (58\%)
(Table~\ref{tab:headline}).

These 14 violations are not concentrated among the classes LTT fails to certify.
Of the 11 classes certified by LTT, 8 (AMI, IMI, IVCD, LVH, NST\_, SARRH, SBRAD,
STTC) violate $\alpha_\mathrm{tier}$ under both uncorrected baselines on the test
fold; the remaining 3 certified classes---IRBBB, LAFB/LPFB, STACH---do not violate
under either baseline.
Conversely, of the 13 classes LTT cannot certify (uncertifiable or trivial),
7 (\texttt{\_AVB}, CLBBB, ISCA, ISCI, ISC\_, RAO/RAE, CRBBB) do not violate
$\alpha_\mathrm{tier}$ under the uncorrected baselines on this fold, while the
remaining 6 (AFIB, LMI, SVARR, LAO/LAE, PACE, RVH) do.
The uncorrected baselines' agreement with the target on these 7 classes is an
outcome of this particular validation/test split rather than a property guaranteed
to hold under resampling, since the baseline procedure carries no statistical
correction for sampling variability between folds.
One of the seven, RAO/RAE, holds at exact equality (test FNR $=\alpha{=}0.10$),
with no margin between the observed value and the target.

Table~\ref{tab:mechanism} reports threshold and test-fold FNR for the 8 classes on
which LTT certifies and both uncorrected baselines violate.
In all 8 rows, the uncorrected threshold exceeds the LTT threshold for both
baselines ($\lambda_\mathrm{TS}>\lambda_\mathrm{LTT}$ and
$\lambda_\mathrm{MC}>\lambda_\mathrm{LTT}$, with no exception across the 16
comparisons).
A higher threshold corresponds to a lower flag rate and, by the definition of FNR
used here, a larger fraction of true positives falling below threshold and being
scored negative---the numeric pattern underlying the violations.
The LTT threshold in each row is set by a procedure that incorporates the sampling
gap between the validation fold, on which $\lambda$ is chosen, and the test fold,
on which FNR is subsequently measured; the uncorrected procedure selects the
threshold at which validation-fold FNR first meets $\alpha_\mathrm{tier}$, with no
explicit allowance for that gap.

Temperature scaling was also evaluated for its effect on aggregate calibration
quality, independent of the per-class violation analysis above.
Fitting $T{=}1.115$ reduced expected calibration error (ECE) from 0.0049 to 0.0023,
a 53\% relative reduction.
This improvement in aggregate calibration is distinct from per-class adherence to a
false-negative-rate target: temperature scaling achieves a substantially lower ECE
while simultaneously violating $\alpha_\mathrm{tier}$ on 14 of 24 classes
(Table~\ref{tab:headline}), indicating that improved aggregate calibration does not
by itself imply per-class control of the false-negative rate.

\subsubsection{Comparison of tiered and uniform-$\alpha$ strategies}

The tiered assignment was compared against three uniform-$\alpha$ strategies
($\alpha{=}0.02$, 0.05, 0.10 applied identically to all classes)
(Figure~\ref{fig:uniform_vs_tiered}; Table~\ref{tab:uniform_comparison}).
Relative to the tiered assignment, the uniform $\alpha{=}0.02$ strategy failed to
certify 9 Benign-tier classes that the tiered assignment certified (efficiency
losses).
The uniform $\alpha{=}0.10$ strategy certified 6 Critical- and Important-tier
classes at a target looser than their assigned tier required (safety failures);
for these classes the published guarantee exceeded the tier's $\alpha$
% CHANGED: "effective FNR bound of 0.054" corrected to "FNR guarantee of 0.10".
% 0.054 (= lambda_c for AFIB under u0.10) is the threshold, not the FNR bound.
% The FNR bound certified by the binomial LTT test at alpha=0.10 is 0.10.
(e.g., AFIB certified at an FNR guarantee of $\alpha{=}0.10$ against a
Critical-tier requirement of $\alpha{=}0.02$).
The two effect sets were disjoint (intersection~$=$~0): no class appeared in both
the efficiency-loss and safety-failure sets.
A further 9 classes were unaffected by the choice of strategy, being uncertifiable
or trivial under all configurations.
In total, 15 of 24 supported classes (63\%) had a different certification outcome
under at least one uniform strategy than under the tiered assignment.

\subsubsection{Discrimination and operating points}

Per-class discrimination, summarised by the area under the receiver operating
characteristic curve (AUROC), ranged from 0.746 to 0.998 across the supported
classes (Figure~\ref{fig:class_auroc}).
Certification outcome and AUROC were not aligned in rank: several high-AUROC classes
were uncertifiable (AFIB, AUROC~$=$~0.983; CLBBB, AUROC~$=$~0.997), while a
lower-AUROC class was certified (IVCD, AUROC~$=$~0.746).
The certified configurations yielded 34 risk-coverage operating points across
16~classes and the three $\alpha$ settings (Figure~\ref{fig:risk_coverage}), with
flag rate and test FNR varying by tier and target as tabulated in
Table~\ref{tab:primary_calibration}.
